\chapter{Conclusion}
\label{ch:conclusion}

A Security Operations Centre rarely fails because an intrusion went undetected. It fails at the step after detection, where a finite number of specialists must decide which of an effectively unbounded stream of alerts deserves their attention. Language models are the first technology capable of reading an event in its context and explaining it in words a person can act on, yet placing one in an operational stream imposes costs in time, trust and data sovereignty that no headline accuracy figure addresses. This thesis asked whether that capability can be obtained without paying its full price, and answered by building SENTINEL---a two-tier architecture that settles cheap decisions before expensive ones, deployed entirely on a local workstation---and then measuring it rather than describing it. Three results carry the claim.

The economics of tiering hold, and hold as a law rather than a single number. Tier 1 and the Machine Learning Gate resolve 90.6\% of a benchmark stream carrying more than three times the realistic attack proportion, and 97.5\% of a stream at a near-realistic base rate, in both cases without a single language model invocation. Because the typical event never reaches the model, median end-to-end latency falls to 0.88\,ms from 17,174.7\,ms under a single-tier baseline. The methodological result matters more than either figure: offload decomposes linearly into a benign-case and an attack-case rate, making it a function of the stream's attack base rate rather than a constant of the system---a property that lets the measurement transfer to another environment, and that a single quoted percentage cannot support.

The security guarantee is structural, and therefore does not decay. Delimited Data Encapsulation treats content inside a per-batch nonce boundary as passive text regardless of what it says, and resisted all 678 adversarial samples in the categories the inexpensive static pre-filter cannot absorb---precisely the categories that filter is blind to. This is a different kind of assurance from a classifier that blocks $x\%$ and leaks the remainder against every novel phrasing. HMAC-SHA256 chaining detects 100\% of the three core tampering modes with no false alarm over an intact control chain, sealing the forensic record with integrity and tamper-evidence.

Hallucination control became a verifiable constraint rather than a hope. Every ATT\&CK identifier the system emits must appear in that batch's retrieved documents; across 1,421 adjudications asserting a technique, zero were unanchored. That zero is not the trivial score of a shield that never acts: it intervened on 9.26\% of batches, and 76 of those interventions withheld a permanent IP block the system could not justify from evidence. The reasoning tier also proved most valuable in a role other than the obvious one---used as a triage router rather than an autonomous confirmer, its deferral queue concentrates 95.0\% of genuine threats into 15.76\% of the volume, a 6.03$\times$ enrichment that removes 84.24\% of the material an analyst must read.

The architecture's limits are reported in full where their evidence lives, in Chapter~\ref{ch:experiments_evaluation}, and three set the immediate agenda: attribution is currently better without the reasoning tier than with it, the agent cannot yet dismiss a false alarm on its own, and the explanation layer is the weakest link, since retrieval admits irrelevant material and only 11.2\% of justifications cite a log value that can be checked back against the record. The corresponding directions are a retrieval layer repaired at the knowledge-base level, calibration for dismissal as well as confirmation, and a justification format that compels citation rather than merely permitting it. Two smaller extensions follow from properties already stated: anchoring hash values periodically outside the system would close the tail-truncation blind spot inherent to log chaining, and asymmetric signatures would upgrade tamper-evidence to genuine non-repudiation.

\vspace{0.3cm}
\begin{flushright}
    \textit{Concluding a rigorous research endeavor, while opening inspiring horizons in securing the digital realm through the power of Agentic Artificial Intelligence.}
\end{flushright}
